\chapter*{Introduzione}
\label{chap:introduzione}
\addcontentsline{toc}{chapter}{Introduzione}
\markboth{Introduzione}{Introduzione}

Negli ultimi due decenni, i progressi dell'apprendimento profondo hanno modificato radicalmente il modo in cui i sistemi artificiali elaborano informazioni visive. Reti neurali convoluzionali prima e architetture basate su Transformer successivamente hanno raggiunto prestazioni elevate in compiti quali classificazione, riconoscimento di azioni, rilevamento di oggetti e comprensione di sequenze video. Una parte considerevole di questa evoluzione si è però sviluppata assumendo implicitamente uno specifico modello di acquisizione del mondo: quello della telecamera convenzionale, che osserva una scena attraverso una successione di immagini complete acquisite a intervalli temporali regolari.

Questo paradigma è estremamente efficace, ma non è neutrale rispetto al fenomeno osservato. La telecamera stabilisce *quando* produrre una nuova misura indipendentemente da quanto la scena sia cambiata. Se una regione dell'immagine rimane immobile, essa viene nuovamente acquisita nel fotogramma successivo; se invece un oggetto si muove rapidamente tra due istanti di campionamento, parte della sua dinamica può essere rappresentata solo indirettamente o può essere degradata da fenomeni quali il *motion blur*. L'informazione visiva viene quindi organizzata secondo il ritmo imposto dal sensore, non secondo quello con cui il mondo evolve.

Le telecamere a eventi adottano un principio differente. Ispirate al funzionamento dei sistemi sensoriali biologici e sviluppate nell'ambito dell'ingegneria neuromorfica, non producono continuamente immagini complete della scena. Ogni elemento fotosensibile reagisce invece alle variazioni locali della luminosità e genera un evento quando il cambiamento osservato supera una determinata soglia. L'uscita del sensore non è pertanto una sequenza regolare di fotogrammi, ma un flusso asincrono di eventi, ciascuno associato a una posizione spaziale, a un istante temporale e al segno della variazione misurata [1]. Questa modalità di acquisizione permette di raggiungere elevata risoluzione temporale, bassa latenza e ampio intervallo dinamico, evitando al tempo stesso di trasmettere ripetutamente gran parte dell'informazione che rimane invariata nella scena [1].

La differenza potrebbe sembrare, a una prima lettura, soltanto una diversa codifica dello stesso contenuto visivo. In realtà modifica la natura del problema computazionale.

In un'immagine convenzionale, il singolo fotogramma costituisce già una rappresentazione spaziale completa del contenuto osservato in un determinato istante. Un algoritmo può quindi analizzarlo come un oggetto sostanzialmente autosufficiente. In un flusso di eventi questa proprietà non esiste. Un singolo evento comunica soltanto che, in una particolare posizione, si è verificata una variazione in un determinato momento. Il significato emerge dalla relazione tra molti eventi distribuiti nello spazio e nel tempo. La traiettoria di un oggetto, la velocità di un gesto, una breve accelerazione, una pausa o la successione di piccoli movimenti non sono necessariamente descritti da un singolo campione: diventano riconoscibili osservando l'evoluzione dell'attività.

Il tempo, di conseguenza, non svolge semplicemente il ruolo di etichetta utilizzata per ordinare misure indipendenti. Esso partecipa alla descrizione del fenomeno.

Si immagini, per esempio, di osservare due movimenti che attraversano le stesse posizioni ma in ordine opposto. Se si ignorasse l'ordine temporale e si sovrapponessero tutte le posizioni visitate, le due osservazioni potrebbero apparire quasi identiche. Eppure rappresenterebbero due movimenti diversi. Analogamente, due sequenze potrebbero generare attività nelle medesime regioni dello spazio, ma una potrebbe consistere in un movimento rapido seguito da una pausa, mentre l'altra potrebbe presentare una progressione lenta e continua. Anche in questo caso la sola distribuzione spaziale dell'attività non sarebbe sufficiente per distinguere completamente i fenomeni.

Questo semplice esempio evidenzia una proprietà fondamentale della visione event-based: **la struttura temporale può contenere informazione discriminativa che scompare quando gli eventi vengono considerati soltanto nel loro insieme**. Un modello efficace non dovrebbe quindi limitarsi a riconoscere *dove* si è verificata attività, ma dovrebbe essere in grado, quando il compito lo richiede, di comprendere anche *quando*, *in quale ordine* e *su quale scala temporale* tale attività si è sviluppata.

Il problema diventa ancora più rilevante nei fenomeni caratterizzati da dinamiche complesse. Un movimento reale non evolve necessariamente secondo una singola scala temporale. Una sequenza può contenere variazioni molto rapide sovrapposte a trasformazioni più lente; brevi burst di eventi possono essere separati da intervalli relativamente inattivi; l'informazione necessaria a prendere una decisione può essere disponibile soltanto dopo aver osservato una parte considerevole della sequenza. Il sistema di elaborazione deve quindi confrontarsi contemporaneamente con dipendenze a breve termine e con relazioni che si estendono per intervalli temporali più lunghi.

Queste caratteristiche pongono una domanda che non esiste nella stessa forma nella classificazione di immagini statiche: **come deve essere progettato un modello affinché sfrutti realmente la struttura temporale del segnale, invece di limitarsi a elaborare una successione di rappresentazioni spaziali?**

Le *Spiking Neural Networks* (SNN) rappresentano una risposta particolarmente interessante a questa domanda. A differenza delle reti neurali artificiali convenzionali, nelle quali l'attività di un neurone viene generalmente rappresentata attraverso valori continui calcolati strato dopo strato, i neuroni spiking possiedono uno stato interno che evolve nel tempo e possono comunicare attraverso eventi discreti, gli *spike* [2]. La risposta di un neurone non dipende quindi soltanto dall'ingresso ricevuto nell'istante corrente, ma anche dallo stato accumulato attraverso gli istanti precedenti. In questo senso la dimensione temporale non viene aggiunta esternamente al modello: appartiene direttamente alla sua dinamica.

La corrispondenza con le telecamere a eventi è immediatamente suggestiva. Da una parte si dispone di un sensore che trasmette informazione sotto forma di variazioni sparse nel tempo; dall'altra, di un paradigma computazionale nel quale l'informazione può essere integrata attraverso stati dinamici e comunicata attraverso eventi discreti. È per questa ragione che le SNN sono frequentemente considerate candidate naturali per l'elaborazione dei dati neuromorfici.

Questa affinità, tuttavia, non deve essere trasformata in un'assunzione.

Il fatto che una rete possieda neuroni dotati di memoria non implica automaticamente che essa abbia imparato a utilizzare in maniera efficace l'ordine temporale presente nei dati. Uno stato neuronale può conservare parte dell'informazione passata e allo stesso tempo essere incapace di rappresentare adeguatamente determinate dipendenze a lungo termine. La rete può inoltre raggiungere buone prestazioni sfruttando regolarità spaziali o statistiche correlate alle classi, senza fare un uso sostanziale della dinamica temporale che si vorrebbe modellare. In altre parole, **essere una rete temporale per costruzione non equivale necessariamente a essere un buon modello del tempo**.

Questa distinzione rappresenta uno dei punti di partenza della presente tesi.

Essa assume particolare importanza perché gran parte della valutazione delle reti neurali, comprese le SNN, viene tradizionalmente riassunta mediante metriche calcolate sull'output finale: accuratezza, F1-score o altre misure di classificazione. Queste metriche sono indispensabili per stabilire se un modello svolga correttamente il compito assegnato, ma descrivono soltanto il risultato della computazione. Due modelli possono ottenere la stessa accuratezza finale e arrivare alla decisione attraverso processi molto diversi: uno può integrare progressivamente informazione lungo l'intera sequenza, mentre l'altro può fare affidamento quasi esclusivamente su caratteristiche disponibili in pochi istanti o su correlazioni che sopravviverebbero anche alterando l'ordine temporale.

La domanda «quanto è accurato il modello?» deve quindi essere accompagnata, quando si studia esplicitamente la capacità temporale, da una seconda domanda: **che uso sta facendo il modello del tempo?**

Parallelamente alla ricerca sulle SNN, un'altra famiglia di architetture ha trasformato profondamente il campo dell'apprendimento automatico: i Transformer. Introdotti originariamente per il trattamento del linguaggio naturale, i Transformer utilizzano meccanismi di *self-attention* per costruire relazioni tra gli elementi di una rappresentazione anche quando questi sono distanti tra loro. Il principio si è dimostrato estremamente efficace anche nella visione artificiale e ha motivato, negli ultimi anni, numerosi tentativi di integrare attenzione e computazione spiking.

Sono così emerse architetture note complessivamente come **Spiking Transformer**. Modelli quali Spikformer hanno mostrato che un meccanismo di self-attention può essere riformulato utilizzando rappresentazioni spiking [3]. Spike-driven Transformer ha ulteriormente spinto questa direzione, cercando di rendere il flusso di calcolo maggiormente compatibile con una computazione guidata dagli spike e riducendo operazioni costose tipiche dei Transformer convenzionali [4]. Più recentemente, QKFormer ha proposto una struttura gerarchica basata su Q-K Attention, nella quale l'importanza dei token o dei canali può essere modellata attraverso rappresentazioni binarie con complessità lineare rispetto alla dimensione considerata [5].

Questi lavori hanno dimostrato che la combinazione tra SNN e Transformer non è soltanto possibile, ma può raggiungere prestazioni competitive mantenendo alcune delle proprietà che rendono interessante il paradigma spiking. Allo stesso tempo, hanno reso più evidente un problema meno immediato.

Nei Transformer convenzionali, l'espressione «attenzione globale» può suggerire l'idea che il modello possieda automaticamente una ricca capacità di relazione anche lungo il tempo. In uno Spiking Transformer applicato a una sequenza, tuttavia, occorre distinguere con precisione **su quale dimensione viene effettivamente calcolata l'attenzione**. In numerose architetture, il meccanismo di attenzione mette in relazione token o canali appartenenti allo stesso intervallo temporale; il collegamento tra un intervallo e quelli precedenti continua invece a dipendere principalmente dallo stato dei neuroni spiking. L'attenzione può dunque essere molto potente nella costruzione delle relazioni spaziali senza realizzare direttamente una corrispondente interazione tra istanti distanti.

La letteratura più recente ha iniziato ad affrontare esplicitamente questo limite. STAtten introduce un meccanismo di attenzione spazio-temporale in grado di elaborare congiuntamente rappresentazioni appartenenti a più istanti, mostrando che l'integrazione delle relazioni temporali può migliorare differenti backbone Spiking Transformer [6].  Il successivo lavoro sul *Multi-Delay Mixer* arriva a una conclusione ancora più esplicita: analizzando il contributo delle informazioni provenienti dagli istanti precedenti, gli autori osservano che diversi meccanismi di attenzione spiking rimangono fortemente concentrati sull'istante corrente e propongono quindi l'integrazione di informazioni ritardate su più scale temporali [7]. 

Si arriva così a un punto fondamentale. La presenza simultanea di neuroni dinamici e di un meccanismo di attenzione non garantisce, da sola, una modellazione temporale sufficientemente ricca. I neuroni spiking forniscono memoria; l'attenzione fornisce potenti interazioni tra rappresentazioni; ma **il modo in cui la capacità del modello viene distribuita tra elaborazione spaziale e temporale rimane una scelta architetturale**.

È proprio questa scelta a diventare critica quando il modello deve essere compatto.

Molte architetture neurali migliorano aumentando progressivamente larghezza, profondità e numero dei parametri. Questo approccio è ragionevole quando il principale obiettivo consiste nel massimizzare le prestazioni e la disponibilità di risorse computazionali non costituisce un vincolo stringente. Diventa però meno soddisfacente quando si vuole comprendere quale componente dell'architettura sia realmente responsabile del miglioramento oppure quando il modello deve essere eseguito su piattaforme con risorse limitate.

Nel caso delle SNN, questa considerazione assume una dimensione ulteriore. Il costo di una rete non dipende soltanto dal numero dei suoi parametri. Devono essere mantenuti gli stati interni dei neuroni; devono essere elaborati gli spike prodotti nel tempo; eventuali buffer temporali occupano memoria; il vantaggio della sparsità esiste soltanto se il sistema di esecuzione può effettivamente evitare le operazioni associate all'attività assente. Un'architettura nominalmente «spiking» può quindi perdere buona parte dei propri vantaggi se introduce operazioni dense, richiede grandi quantità di stato o genera attività eccessivamente elevata.

Per questo motivo la prospettiva hardware non viene considerata in questa tesi come una fase completamente separata dalla progettazione del modello. L'obiettivo di un futuro deployment su FPGA impone di valutare fin dall'inizio non soltanto quanti parametri possieda un'architettura, ma anche come essa utilizzi memoria e calcolo. Le FPGA offrono infatti parallelismo configurabile, memoria locale e possibilità di realizzare datapath specializzati, ma dispongono di risorse finite in termini di LUT, flip-flop, DSP e BRAM. La letteratura sull'implementazione di SNN su FPGA evidenzia inoltre che lo sfruttamento della sparsità richiede scelte architetturali specifiche: evitare computazioni sugli ingressi inattivi è vantaggioso solo quando l'overhead necessario a gestire il comportamento sparso rimane inferiore al costo delle operazioni risparmiate [8]. 

Emergono quindi tre esigenze che devono essere considerate congiuntamente.

Il modello deve essere capace di utilizzare la struttura temporale degli event stream; deve farlo in un regime fortemente compatto, nel quale non sia possibile risolvere ogni limite aumentando indiscriminatamente la capacità; infine, le modifiche introdotte devono rimanere coerenti con una futura realizzazione hardware efficiente.

Queste esigenze sono tra loro strettamente connesse. Aggiungere capacità temporale può aumentare memoria e operazioni. Ridurre eccessivamente la dimensione del modello può compromettere la qualità delle rappresentazioni. Massimizzare la sparsità senza considerare quale informazione venga eliminata può produrre una rete efficiente ma poco accurata. La progettazione diventa pertanto un problema di **allocazione della capacità**: non soltanto quante risorse assegnare al modello, ma soprattutto *dove* impiegarle e *quale funzione* debbano svolgere.

La presente tesi nasce all'interno di questo spazio progettuale.

Il punto di partenza è un Transformer spiking compatto derivato da QKFormer. L'obiettivo non è costruire progressivamente un modello sempre più grande né cercare una combinazione di iperparametri capace di massimizzare una singola metrica. Il lavoro intende piuttosto studiare in modo controllato quali componenti siano effettivamente necessarie per elaborare stream neuromorfici temporalmente complessi, quali limiti emergano quando la capacità complessiva è vincolata e se una diversa distribuzione delle risorse possa migliorare la modellazione temporale senza compromettere l'efficienza.

Questa impostazione comporta anche una precisa scelta metodologica. Prima di intervenire sull'architettura occorre comprendere il comportamento della baseline. Prima di interpretare un aumento dell'accuratezza come evidenza di una migliore modellazione temporale occorre verificare quali informazioni il modello utilizzi. Prima di attribuire un vantaggio a un nuovo modulo occorre confrontarlo con controlli che permettano di separare l'effetto della struttura dall'effetto del semplice aumento della capacità.

Il percorso seguito nel lavoro procede quindi dalla **misura alla diagnosi e dalla diagnosi alla progettazione**, anziché dalla progettazione a una successione di tentativi empirici. Una prima parte dell'analisi indaga se e quanto un modello spiking apparentemente competente sfrutti realmente l'ordine e l'evoluzione temporale degli input. Su questa base viene poi costruita e stabilizzata una baseline compatta; vengono individuati i principali limiti attraverso esperimenti controllati; infine vengono studiate modifiche architetturali che agiscono su aspetti distinti della computazione temporale e dell'allocazione delle risorse. Solo dopo questa fase la valutazione può essere estesa alla robustezza statistica, alla capacità di generalizzazione e alle implicazioni hardware.

Il problema centrale affrontato dalla tesi può essere quindi espresso in forma ancora preliminare nel modo seguente: **come progettare uno Spiking Transformer compatto che non si limiti a elaborare dati distribuiti nel tempo, ma che utilizzi efficacemente la loro struttura temporale, mantenendo al contempo un profilo computazionale compatibile con un'implementazione hardware efficiente?**

Le sezioni successive costruiranno progressivamente gli strumenti necessari per rendere questa domanda precisa. Sarà dapprima necessario comprendere che cosa distingua l'acquisizione event-based dalla visione convenzionale e quale informazione venga effettivamente trasportata da un event stream. Verranno quindi introdotte le dinamiche fondamentali delle reti neurali spiking e il ruolo dello stato neuronale nella memoria temporale. Su questa base sarà possibile esaminare l'integrazione tra SNN e Transformer, fino a identificare con maggiore precisione il limite che motiva il lavoro: la necessità di conciliare **modellazione temporale, compattezza architetturale ed efficienza hardware** all'interno dello stesso sistema.
